\section{Introduction}
\label{sec:intro}

This paper provides causal evidence that uninsured depositors
discipline community banks.  When mandatory adoption of the Current
Expected Credit Loss (CECL) standard forced more than 4,000
U.S.\ community banks to disclose lifetime expected credit losses on
their existing loan portfolios in 2023, banks with large disclosed
charges lost 0.39 percentage points of assets in uninsured time
deposits, about 7 percent of the sample mean, and gained a similar
amount of insured time deposits.  Whether such discipline exists is
a central question in banking.  Theory predicts it should
\citep{calomiris1991role, diamond2001liquidity}, and riskier banks pay
higher deposit rates and experience larger uninsured outflows
\citep{park1995market, martinez2001depositors, maechler2006dynamic,
iyer2016tale, egan2017deposit, chen2024liquidity}, but whether these
patterns reflect a causal disciplinary response or equilibrium
correlation remains unsettled.  The question is pressing: the Main
Street Depositor Protection Act (S.~2999, reintroduced March 2026)
would extend coverage to \$10~million per account, and early state
deposit insurance programs that eliminated market discipline enabled
excessive risk-taking \citep{calomiris2019stealing}.  Community banks
are the right laboratory: institutions below \$10~billion in assets
carry no credible government backstop, so the monitoring incentive is
intact \citep{o1990deposit, flannery1996evidence}.

Clean identification requires resolving three problems.  Simultaneity:
a bank's risk profile, deposit composition, and depositor behavior
co-evolve, so regressions of flows on risk measures capture equilibrium
sorting rather than isolated disciplinary responses.  Accounting
opacity: under the incurred-loss model, banks could delay recognition
and systematically understate risk in disclosures
\citep{bushman2012accounting, bushman2015delayed}.  Aggregate stress:
periods with the greatest variation in bank distress are precisely those
in which government interventions and depositors' own liquidity needs
contaminate the response \citep{bennett2015market}.  I address all
three by exploiting mandatory CECL adoption for community banks in 2023
as an information shock whose cross-sectional intensity was fixed
before the event.  CECL requires banks to recognize lifetime expected
credit losses on existing loan portfolios at adoption, producing a
one-time adjustment to retained earnings whose magnitude reflects loan
vintages assembled years earlier.  Banks knew the standard was coming,
but the size of the Day-One charge, the treatment variation I use, was
pinned down by the composition and seasoning of the existing loan book
rather than by contemporaneous deposit-market conditions.
\citet{kim2026current} and \citet{gee2022decision} confirm the Day-One
adjustment is informative; \citet{nier2006market} and
\citet{chen2022bank} establish that disclosure improvements enable
depositors to act on credit-risk signals.

The central result, estimated on quarterly data from 2016 to 2025, is
the level difference-in-differences estimate above: after the
disclosure, high-CECL banks lost 0.394 percentage points of assets in
uninsured time deposits and gained 0.712 percentage points of insured
time deposits, with continuous-treatment estimates scaling in
proportion to the disclosed charge.  Event studies over a five-year
pre-period show flat pre-adoption trajectories for the uninsured
series and a divergence that emerged at adoption and persisted.  As a
very conservative robustness check, a triple-difference specification
with bank-by-quarter fixed effects, which compares uninsured to
insured time deposits within the same bank and quarter and absorbs
every bank-level shock, confirms the composition shift: each
percentage point of CECL exposure moves uninsured time deposits by
0.18 percentage points of assets relative to insured time deposits.
The response is confined to the margin where discipline is feasible:
time deposits reprice only at maturity, and the rollover decision is
the depositor's to make.

The timing of the response separates the two actors in the market.
Banks knew their Day-One charges before depositors could: supervisors
and auditors expected banks to run CECL in parallel with the
incurred-loss model through 2022, and the anticipatory behavior I
document is itself direct evidence that management had a quantified
estimate before the public filing.  Bank-controlled margins move
first.  Insured brokered deposits at high-CECL banks begin
accumulating three quarters before adoption, and the size of that
pre-adoption ramp is predicted by the still-undisclosed future charge;
banks with larger coming charges were also six times more likely to
elect the regulatory capital phase-in, a direct revelation that they
expected a material hit.  Depositor-facing outcomes move only when the
information becomes public: uninsured time deposits are flat through
every pre-adoption quarter and break at disclosure.

The reallocation carries real costs.  High-CECL banks paid
approximately 8--9 basis points more annually in interest expense, ROA
fell by roughly 15 basis points, and NIM compressed by approximately 9
basis points.  Two complementary rate tests show where the cost arises
and where it does not.  Implied rates from Call Report interest
expense show high-CECL banks did not pay more on the uninsured
certificates they retained, and weekly posted deposit rates from
RateWatch show they never raised retail CD or savings rates.  The
funding-cost increase reflects the substitution itself: maturing
uninsured certificates replaced with insured brokered and reciprocal
money raised at wholesale rates.  Discipline takes the form of a
quantity response: depositors leave rather than negotiate, and the
bank pays for their departure through the higher cost of the insured
wholesale funding that replaces them.

Community banks adopted CECL in the same quarter that the failures of
Silicon Valley Bank and Signature Bank put bank risk on every
depositor's mind, so I confront the stress confound directly.  Four
pieces of evidence separate the disclosure response from the 2023
stress.  First, the disclosed adjustment is essentially uncorrelated
with the balance-sheet characteristics that governed run exposure in
2023: unrealized securities losses, uninsured deposit shares, and the
other vulnerability measures jointly explain less than one percent of
its cross-sectional variation, and the uninsured share enters with the
wrong sign for a stress story.  Second, the triple-difference
comparison nets out any bank-level exposure to the stress by
construction.  Third, high-CECL banks' posted deposit rates show no
differential movement during the stress weeks themselves.  Fourth, the
bank-side response began in 2022, before the failures.  Separately,
the large SEC-filing banks that adopted CECL in 2020 validate the
signal content of the disclosure: their market capitalization and CDS
spreads moved sharply at adoption in proportion to the charge.  I
exclude that cohort from the deposit tests because it is small
(roughly one hundred publicly traded holding companies) and because
the COVID-19 liquidity programs flooded banks with deposits in
2020, swamping any disciplinary response.

This paper contributes to three literatures.  First, on depositor and
market discipline, prior evidence documents higher rates and larger
outflows at riskier banks \citep{park1995market, egan2017deposit,
flannery1996evidence, nier2006market}, but causal interpretation has
been elusive: government guarantees attenuate estimates in the crisis
episodes that generate the most variation \citep{bennett2015market,
berger2015depositors, jacewitz2018deposit}, and cross-sectional
comparisons cannot separate information responses from equilibrium
sorting.  Recent work on the 2023 episode shows depositors ran on banks
with high uninsured shares and unrealized losses
\citep{caglio2024depositors}; I exploit orthogonal variation
\emph{within} that high-attention period, uncorrelated with the
balance-sheet characteristics that governed run exposure.  Second, on
CECL and mandatory credit-risk disclosure, \citet{bushman2012accounting}
and \citet{bushman2015delayed} document that delayed loss recognition
under the incurred-loss model attenuated the informativeness of bank
statements and, by extension, the market's ability to price bank risk.
Prior CECL work focuses on the supply side, lending
\citep{granja2025current} and provisioning \citep{kim2026current}, and
on information users other than depositors, particularly analysts
\citep{gee2022decision} and equity investors.  I complement this
literature by tracing the disclosure to depositor behavior and
downstream to funding costs and profitability, closing a channel that
the analyst-and-equity strand does not observe.  Third, on deposit
insurance design, broader coverage reduces the incentive to monitor
\citep{demirgucc2004market, flannery2019market}; the substitution I
document toward insured brokered and reciprocal deposits
\citep{kim2024insurance} shows that much of the response takes the form
of repurchasing insurance through market arrangements rather than
exiting the banking system.
